% Source material: docs/seminar/audit.md — implementation details
% See also: src/csttool/ source code for pipeline architecture

\chapter{Methods}
\label{ch:methods}

\section{Pipeline Overview}
\label{sec:methods_pipeline}

TODO: Brief intro to the standard pipeline. Add mermaid diagram or flowchart. Mention every module can run individually, as a full run, or multiple runs via batch.

\section{Data Import}
\label{sec:methods_dataimport}

TODO: Explain import pipeline, how DICOMs and NIfTIs are handled, gradient table construction, validation of input data structure.

\section{Preprocessing}
\label{sec:methods_preprocessing}

TODO: Explain all parameters and their CLI calls in a table. Explain how this package is underdeveloped to keep it Pythonic and leaner.

\section{Diffusion Tensor Fitting}
\label{sec:methods_diffusiontensorfitting}

TODO: Tie in to \ref{subsec:background_dti}; describe the use of DIPY's TensorModel, the eigenvalue decomposition, derived scalar maps (FA, MD, AD, RD), white matter mask generation, NaN/Inf handling. 

Question: FA map as redundancy if anat T1 not found for future pipeline?

\section{Whole Brain Tractography}
\label{sec:methods_wholebraintractography}

TODO: Tie in to \ref{subsec:background_tractography}. Points to touch on:

- Direction field estimation:
- CSA-ODF model (Constant Solid Angle)
- Spherical harmonic order (default 6, auto-reduced if insufficient gradients)

Seeding:
- Seed mask: voxels with FA $\geq$ 0.2
- Seed density: configurable (default 1 per voxel)

Propagation:
- Deterministic local tracking (DIPY's LocalTracking)
- Step size: 0.5 mm (Eq. 2.10 from your background)
- Fixed random seed (default 42) for reproducibility

Stopping criterion:
- FA threshold-based (ThresholdStoppingCriterion)
- Optional brain mask boundary

Justify deterministic over probabilistic: reproducibility, interpretability, sufficient for CST (low crossing-fiber complexity).


\section{Atlas-Based CST Extraction}
\label{sec:methods_atlasbasedextraction}

TODO: Maps to \ref{subsec:background_atlasbasedextraction}

\subsection{Image Registration}
\label{subsec:methods_imageregistration}

TODO: How the atlas-to-subject approach described in \ref{subsubsec:background_registration} is implemented:

1. Affine registration: Translation → Rigid → Affine (Mutual Information metric)
2. Diffeomorphic registration: SyN (Symmetric Normalization, Cross-Correlation metric)
3. Target: MNI152 template (1mm, skull-stripped, bundled with csttool)


\subsection{Atlas Warping}
\label{subsec:methods_atlaswarping}


- Harvard-Oxford atlases (FSL, via templateflow):
	- Cortical: 96 regions
	- Subcortical: 21 regions
- Inverse deformation applied to warp labels MNI → subject native space


\subsection{ROI Definition}
\label{subsec:methods_roidefinition}

TODO: Which ROIs are defined, from which atlas, label number, purpose. Maybe tie in to \ref{sec:background_anatomicalbackground}

\subsection{Streamline Filtering}
\label{subsec:methods_streamlinefiltering}

TODO: Tie into \ref{subsubsec:background_filteringstrategies}. Explain this:

  1. Endpoint filtering (strict): Requires one endpoint in brainstem, other in motor cortex. Bidirectional checking.
  2. Passthrough filtering (default, permissive): Streamlines must traverse both ROIs; endpoints flexible.

Length filtering: 20–200 mm (anatomically plausible CST range).


\section{Tract-Level Metrics}
\label{sec:methods_metrics}

\subsection{Unilateral Analysis}
\label{subsec:methods_unilateral}

Morphological:
- Streamline count
- Length statistics (mean, std, min, max)
- Tract volume (unique voxels x voxel size, mm3)

Microstructural:
- FA, MD, RD, AD sampled along streamlines
- Statistics: mean, std, median, min, max
- 20-point along-tract profiles


\subsection{Bilateral Analysis}
\label{subsec:methods_bilateral}

Laterality Index:

\begin{equation}
	LI = \frac{L-R}{L+R}
\end{equation}

Interpretation thresholds: What makes the result symmetric, mildly asymmetric, moderately asymmetric or strongly asymmetric. Find sources maybe?

Computed for: volume, streamline count, length, FA, MD, RD, AD


\section{Implementation Details}
\label{sec:methods_details}


- Language: Python $\geq$ 3.10
- Core libraries: DIPY (tractography, tensor fitting), nibabel (NIfTI I/O), nilearn (atlas handling), scipy (morphology). MAKE SURE TO CITE THEM ALL.
- CLI architecture: Modular commands, each stage independent
- Reproducibility: Provenance tracking, optional deterministic seeding
- More?

